\section{The Verification Ledger}
\label{sec:ledger}

% 3-4 pages. THE core section. Write this first, or second after Section 2.
% Normative source: implementation.md section 2 (semantics) and section 11
% (executable contract). Every claim here must trace to one of those.

\subsection{Definitions}
% Snapshot \snapshot: immutable content identity of the workspace (git tree
%   hash; submodule and LFS pointers are tree content).
% Action: a step that may mutate the workspace. Completing it yields a new
%   snapshot or leaves the snapshot unchanged. An action has run-state, never
%   validity -- it is not "green forever".
% Gate g: a deterministic predicate over a snapshot, producing pass/fail plus
%   evidence. A gate never mutates the workspace.
% Driver: the one action the workflow designates as the agent. Runs once at
%   task start, and again whenever a failed gate needs a fix that no cheaper
%   fix action resolved.
% Fix action: an ordinary action class referenced by a gate class; tried once
%   per snapshot before a driver run is spent.
% Gate key \gatekey = hash(snapshot, resolved gate specification): changing
%   the verifier invalidates its old results, exactly as in a build system.
% Ledger \ledger: partial function (g, \gatekey) -> {pass, fail(e)}.
%   Kubernetes materialises it as a BOUNDED view (current result and verified
%   snapshot per gate, explicit counters), not an append-only log (Section 8.4).
% Formalise: one \begin{definition} block per item; Appendix B carries the
% formal statements, keep this subsection readable.
% The sentence to land: actions mutate state; gates assert properties of state.
% Everything in the model follows from keeping those two apart.

\subsection{Path-scoped invalidation}
% Each gate declares invalidatedBy: a set of path globs.
% On an edit producing h' from h, gate g retains its result iff
%   diff(h, h') INTERSECT invalidatedBy(g) = empty set.
% THE DEFAULT IS CLOSED: an undeclared invalidatedBy means ["**"] -- any change
% invalidates. Carry-forward is an explicit optimisation the author opts into.
% Argue the asymmetry: a false carry-forward is a correctness and security bug;
% an unnecessary re-run is merely expensive. (Section 9.3 shows how the open
% default would have neutralised the injected integrity gate.)
% Failed results carry forward under the same rule as passed ones; a truncated
% or contradictory change report is treated as ["**"].
% Two gate scopes: tree-scoped gates key on the snapshot alone; diff-scoped
% gates (the suppression meta-gate, Section 9) key on (base revision, snapshot)
% [spec].
% Corollary worth stating: a rebase is just another tree change. It invalidates
% diff-scoped gates and preserves tree-scoped ones whose paths it misses. This
% is what makes high monorepo concurrency affordable -- and it is why the base
% revision must NOT be folded into snapshot identity itself.

\subsection{\fold and \decide as pure functions}
% Two functions, no Kubernetes dependency, no I/O, no clock:
%   \fold(view, execution result) -> view: adopts the snapshot, applies
%     invalidation, updates counters, detects oscillation.
%   \decide(plan, view, budgets) -> RunAction | RunGate | Escalate | Done.
% Include the pseudocode (implementation.md section 11.3). Points to make:
%   - cost-ordered, fail-fast: steps walk in dependency order; a cheap failing
%     gate short-circuits before an expensive one runs;
%   - fix action before driver: remediation escalates in cost;
%   - budgets guard REMEDIATION, never verification -- the last budgeted fix is
%     always observed before any escalation;
%   - blocking is honoured: a non-blocking failure is recorded and never blocks
%     the fixed point;
%   - the driver is a step: a task whose base tree already passes every gate
%     still runs it. Verifying the base tree is never success.
% Unit-testable against a synthetic ledger. The implementation's test matrix
% (implementation.md section 11.7) is the evidence: cite case counts and the
% property suite, not adjectives.

\subsection{Termination}
% Fixed point: every blocking gate green against a SINGLE snapshot AND the
%   driver completed at least once.
% Oscillation: a newly produced snapshot equals a non-parent ancestor in the
%   lineage => the loop is cycling; break immediately. A no-op action (the
%   snapshot is unchanged) is NOT oscillation.
% A lifetime cap on driver runs plus a per-gate cap on fix attempts replace
%   the sliding-window budget of earlier drafts: with one driver the window
%   bought nothing but complexity.
% Global budgets: tokens; ACTIVE wall clock (accumulates from provisioning,
%   pauses while suspended); consecutive infrastructure failures, escalating
%   only where no fallback exists (driver, blocking gates) -- an exhausted fix
%   action degrades to the driver, an exhausted non-blocking gate is tolerated.
% Diff-size ratchet (monotonic growth with gates red => flailing, revert): an
%   early-exit heuristic layered on termination, not part of it [spec].
% \begin{proposition}: bounded termination with the explicit bound of
%   implementation.md section 11.5 -- driver and fix runs bound snapshots,
%   snapshots bound gate runs, consecutive-reset bounds errors. Proof sketch in
%   Appendix B.

\subsection{Comparison to one-shot execution}
% Table from Section 3.2, expanded. This is where the reader should conclude
% that a DAG engine cannot express this without becoming a ledger.
% Add the row a pipeline cannot have: "gate result becomes unknown because the
% tree changed underneath it" (line 17 of the reference trace).
